%\documentclass[11pt,a4paper,twoside]{article}

\usepackage[utf8]{inputenc}

\usepackage[T1,T2A]{fontenc}

\usepackage[english.ancient,russian.ancient]{babel}

\usepackage{graphicx}

\usepackage{array}

\usepackage[doi]{samgtu-bib}

\usepackage{samgtu}

\usepackage{samgtu-name}

\usepackage{mathtools}

\mathtoolsset{showonlyrefs}

\usepackage{desclist}

\usepackage{euscript}

\usepackage{mathrsfs} 

\usepackage{multicol}

\usepackage{mathrsfs}

\usepackage{cite}

\usepackage{cmap}

\usepackage{qrcode}

\allowdisplaybreaks[4]

\usepackage[nodayofweek]{datetime}

\usepackage[thinlines]{easybmat}



\renewcommand{\JournalYear}{2018}

\renewcommand{\JournalVolume}{22}

\renewcommand{\JournalNumber}{4}



\newdate{JournalDateSign}{29}{12}{2018}% Дата подписания Вестника в печать

\renewcommand{\JournalDateSign}{\displaydate{JournalDateSign}}

\newdate{JournalDatePrint}{16}{01}{2019}% Дата выхода Вестника из печати

\renewcommand{\JournalDatePrint}{\displaydate{JournalDatePrint}}

%\renewcommand{\JournalUPL}{16.56} % Усл. печ. л.

%\renewcommand{\JournalUIL}{16.53} % Уч.-изд. л.

\renewcommand{\JournalUPL}{15.24} % Усл. печ. л.

\renewcommand{\JournalUIL}{15.21} % Уч.-изд. л.

\renewcommand{\JournalRegNumber}{220/18} % Регистрационный номер

\renewcommand{\JournalOrderNumber}{2} % Номер заказа



%\renewcommand{\JournalRMonth}{Март} % Месяц выхода Вестника

%\renewcommand{\JournalEMonth}{March} % Месяц выхода Вестника

%\renewcommand{\JournalRMonth}{Июнь} % Месяц выхода Вестника

%\renewcommand{\JournalEMonth}{June} % Месяц выхода Вестника

%\renewcommand{\JournalRMonth}{Сентябрь} % Месяц выхода Вестника

%\renewcommand{\JournalEMonth}{September} % Месяц выхода Вестника

\renewcommand{\JournalRMonth}{Декабрь} % Месяц выхода Вестника

\renewcommand{\JournalEMonth}{December} % Месяц выхода Вестника



\newcommand{\totop}[1]{\raisebox{6pt}[1pt][2pt]{#1}}

\newcommand{\tottop}[1]{\raisebox{12pt}[1pt][2pt]{#1}} 

\newcommand{\totttop}[1]{\raisebox{18pt}[1pt][2pt]{#1}} 

\newcommand{\tottttop}[1]{\raisebox{24pt}[1pt][2pt]{#1}} 

\newcommand{\totttttop}[1]{\raisebox{30pt}[1pt][2pt]{#1}} 

\newcommand{\tobot}[1]{\raisebox{-6pt}[1pt][2pt]{#1}} 

\newcommand{\tobbot}[1]{\raisebox{-12pt}[1pt][2pt]{#1}} 

\newcommand{\tobbbot}[1]{\raisebox{-18pt}[1pt][2pt]{#1}}



\graphicspath{{./00PIC/}}
%\usepackage{samgtu-name}
%
\vsgtu{1564} %registration number
\FirstPage{773}
\LastPage{785}
%
\ShortCommunication
%
%\pdfcrop
%\draft
%
%
%\begin{document}

\renewcommand{\baselinestretch}{0.885}


\hfill \qrcode[hyperlink,height=.8in]{http://doi.org/10.14498/vsgtu1564}
\vspace{-.8in}


\udc{539.3}
\msc{74F, 74K20}

\rutitle{Вынужденные осесимметричные колебания круглых~многослойных биморфных пластин}

\entitle{Forced axisymmetric oscillations of circular multilayer bimorph  plates}

\ruauthor{Д.~А.~Шляхин, О.~В.~Ратманова}{Ш\,л\,я\,х\,и\,н~Д.~А., Р\,а\,т\,м\,а\,н\,о\,в\,а~О.~В.}
\enauthor{D.~A.~Shlyakhin, O.~V.~Ratmanova}{S\,h\,l\,y\,a\,k\,h\,i\,n~D.~A., R\,a\,t\,m\,a\,n\,o\,v\,a~O.~V.}

\ruaffil{\rusamgtu.}
\enaffil{\ensamgtu.}



\ruauthorsdetails{2}{% Количество авторов
\fullauthor{\rupid{52312}{Дмитрий Аверкиевич Шляхин}}
\CAuthor % Автор-корреспондент
\orcid{0000-0003-0926-7388} % ORCID ID автора 
\regalia{доктор технических наук, доцент} 
\position{заведующий кафедрой} 
\dept{каф. сопротивления материалов и строительной механики} 
\email{d-612-mit2009@yandex.ru}
\fullauthor{\rupid{131941}{Олеся Викторовна Ратманова}}
%\CAuthor % Автор-корреспондент
\orcid{0000-0001-7591-4921}
\position[n]{ассистент} 
\dept{каф. сопротивления материалов и строительной механики} 
\email[b]{olesya654@yandex.ru}
}



\enauthorsdetails{2}{
\fullauthor{\enpid{52312}{Dmitry A. Shlyakhin}}
\CAuthor % Автор-корреспондент
\orcid{0000-0003-0926-7388} % ORCID ID автора 
\regalia{Dr.  Techn. Sci.} 
\position{Head of Dept.} 
\dept{Dept. of Strength of Materials and Structural Mechanics}
\email[v]{d-612-mit2009@yandex.ru}
\fullauthor{\enpid{131941}{Olesya V. Ratmanova}}
%\CAuthor % Автор-корреспондент
\orcid{0000-0001-7591-4921} 
\position[n]{Assistant} 
\dept{Dept. of Strength of Materials and Structural Mechanics}
\email[v]{olesya654@yandex.ru}
}



\rukeywords{биморфная пластина, электроупругость, 
нестационарная нагрузка, интегральные преобразования}

\enkeywords{bimorph plate, electroelasticity, nonstationary loading, integral transformations}

\ruabstract{Представлена методика расчета круглых сплошных многослойных биморфных 
пластин и получены новые замкнутые решения осесимметричных динамических 
задач прямого и обратного пьезоэффектов. В~общем случае при исследовании 
электроупругого (пьезокерамического) и~упругого слоев математическая 
формулировка рассматриваемых задач включает уравнения движения и~Максвелла 
в~пространственной постановке относительно компонент вектора перемещений 
и~потенциала электрического поля, а~также соответствующие начально"=краевые 
условия. Рассмотрены случаи шарнирного и~жесткого закрепления внешнего 
контура конструкции. Для исследования связанных линейных задач применяется 
математический аппарат в виде метода конечных интегральных преобразований 
Фурье"--~Бесселя и обобщенного интегрального преобразования (КИП). При этом на 
каждом этапе исследования использовалась процедура приведения граничных 
условий к виду, позволяющему выполнять соответствующее преобразование.

Построенные расчетные соотношения позволяют обосновать конст\-руктивные 
решения многослойных пьезокерамических преобразователей, а~именно 
подобрать геометрические размеры и физические характеристики используемых 
материалов, определить размеры разрезных круговых электродов, позволяющие 
наиболее эффективно преобразовать внешнее электрическое воздействия 
в~механические колебания при различной частоте. Кроме этого, появляется 
возможность проанализировать напряженно"=деформированное состояние, характер 
изменения электрического поля, а~также частотный спектр собственных 
осесимметричных колебаний рассматриваемых систем.}

\enabstract{A method for calculating circular multilayer bimorph plates is presented and new analytical solutions of axisymmetric dynamic problems of direct and inverse piezoelectric effects are obtained. The cases of hinged and rigid fixing of the outer contour of the structure are considered. To investigate related linear problems, a mathematical apparatus is used in the form of a method of finite integral transformations. The constructed calculated relationships allow us to substantiate the constructive solutions of piezoceramic transducers. 

Built design ratio allows to prove the constructive solutions of multilayer piezoelectric ceramic transducers, namely, to choose the geometrical dimensions and physical characteristics of the materials used, define the dimensions of a split circular electrode, allowing most effectively to convert the external electrical stimulation into mechanical vibrations at various frequency. In addition, it is possible to perform stress-strain state, the nature of the change of the electric field and frequency range of the axisymmetric vibrations of the considered systems.
}



\newdate{shlyakina}{14}{10}{2017}
\date{\displaydate{shlyakina}}
\newdate{shlyakinb}{09}{12}{2017}
\revisiondate{\displaydate{shlyakinb}}
\newdate{shlyakinc}{18}{12}{2017}
\accepteddate{\displaydate{shlyakinc}}

\newdate{shlyakine}{29}{12}{2017}
\renewcommand{\JournalDataOnline}{\displaydate{shlyakine}}

\vspace{5mm}

%\ForPeerReview
\makerutitle



\Section[n]{Введение} В различных электроакустических, оптических и 
пневматических приборах применяются пьезокерамические преобразователи в виде 
тонких биморфных пластин~[\citen{shlya:bib1}--\citen{shlya:bib3}]. Данные конструкции представляют собой 
симметричные или асимметричные многослойные системы, состоящие из 
электроупругих и упругих тонких пластин. 
Их широкое применение объясняется 
высокими эксплуатационными параметрами, низкой стоимостью и простотой в 
управлении. 
Изгибные колебаний в рассматриваемых конструкциях создаются с~помощью электрического напряжения, приложенного к электродированным поверхностям пьезокерамических пластин (явление обратного пьезоэффекта), или 
в случае действия механической нагрузки (прямой пьезоэффект).

Для описания работы и повышения функциональных возможностей биморфных 
пластин возникает необходимость в разработке математических моделей и 
построении общих решений, позволяющих описать сложную картину взаимодействия 
полей напряжений различной физической природы~\cite{shlya:bib4}. 
В~настоящее время для  реализации данной проблемы, как правило, применяются прикладные теории для 
тонких пластин в условиях установившегося режима вынужденных колебаний при 
использовании различных допущений о характере распределения напряженности 
электрического поля в пьезокерамических пластинах~[\citen{shlya:bib5}--\citen{shlya:bib11}]. Полученные 
приближенные расчетные соотношения используются в дальнейшем для расчета и 
конструирования приборов различного назначения~[\citen{shlya:bib12}--\citen{shlya:bib15}]. 

Замкнутые решения динамических начально"=краевых задач электроупругости в 
пространственной постановке при аксиальной поляризации пьезокерамического 
материала получены в работах~[\citen{shlya:bib17}--\citen{shlya:bib19}] только для однородного диска 
(цилиндра). В~\cite{shlya:bib17} рассматривался незакрепленный цилиндр при учете смешанных 
краевых условий на торцевых поверхностях, а в~\cite{shlya:bib18} исследовалось 
напряженно"=деформированное состояние жестко закрепленной толстой пластины в 
случае действия осесимметричной механической нагрузки. Здесь также можно 
отметить одну из работ, посвященную определению спектра частот собственных 
осесимметричных и неосесимметричных колебаний~\cite{shlya:bib19} пьезокерамических  
цилиндров.

Таким образом, в настоящее время отсутствуют замкнутые решения динамических 
задач электроупругости для многослойных конструкций в пространственной 
постановке. Для решения этого вопроса в данной работе представлена методика 
расчета круглых многослойных шарнирно и жестко закрепленных биморфных 
сплошных пластин в случае действия электромеханической нагрузки. Построенные 
расчетные соотношения позволяют дать качественную и количественную оценку 
связанности электрических и механических полей напряжений в многослойных 
конструкциях. Частные решения данного подхода приведены в публикациях 
авторов~\cite{shlya:bib21,shlya:bib22}.

\smallskip

\Section{Постановка задачи} В общем случае круглая биморфная конструкция, 
занимающая в~цилиндрической системе координат область $\Omega= \big\{ (r,\theta ,z) :  0\le r\le b,0\le \theta \le 2\pi 
, 0\le z\le h^\ast\big\}$, представляет многослойную систему 
в~виде жестко соединенных между собой электроупругих (пьезокерамических) и~упругих тонких пластин. 
Для определенности принимаем, что она состоит из 
металлической подложки толщиной $h_2^\ast$ и двух пьезокерамических пластин 
высотой $h_1^\ast$, выполненных из материала гексагональной системы класса 
6mm ($h^\ast =2h_1^\ast +h_2^\ast$). Продольно"=поперечные осесимметричные 
колебания индуцируются за счет подведения к торцевым электродированным 
поверхностям пьезопластин с различным направлением вектора аксиальной 
поляризации электрического напряжения $V^\ast (r_\ast,t_\ast )$ (рис.~сверху) или действии на лицевых плоскостях конструкции механических нагрузки в виде нормальных напряжений $q^\ast  (r_\ast,t_\ast )$ (рис.~снизу). 
В дальнейшем все рассматриваемые величины будут обозначаться с символом $\ast$, а безразмерные "--- без символа.
Рассматривается жесткое и шарнирное закрепления внешнего контура  электроупругой системы (см.~рис.).

\begin{figure}[h!]\centering\small
\includegraphics[scale=1.05]{shlyakhin1}\\
\vspace{5mm}
\includegraphics[scale=1.05]{shlyakhin2}\\
\caption*{Расчетная схема биморфных пластин [Calculation scheme of bimorph plates]}

\vspace{-5mm}
\end{figure}

Математическая формулировка рассматриваемых осесимметричных начально"=краевых задач в безразмерной форме в~цилиндрической системе координат  включает следующие соотношения~\cite{shlya:bib4}:
\begin{itemize}
\item дифференциальные уравнения движения и Максвелла относительно безразмерных компонент 
вектора перемещений $U(r,z,t)$, $W(r,z,t)$, а 
также безразмерного потенциала электрического поля $\varphi (r,z,t)$:
\begin{equation}
\label{shlya:eq1}
\begin{array}{l}
\displaystyle 
\frac{\partial }{\partial r}\nabla U+\frac{C_{55}^{( s )} }{C_{11}^{( s )} }
\frac{\partial ^2U}{\partial z^2}+
\frac{ C_{13}^{( s )} +C_{55}^{( s )} }{C_{11}^{( s )} }
\frac{\partial ^2W}{\partial r\partial z}+
\\[2mm]
\displaystyle \hspace{3cm}
+\frac{ e_{31} +e_{15} }{e_{33}}\frac{\partial ^2\varphi }{\partial r\partial z}-
\Phi^{( s )}\frac{\partial ^2U}{\partial t^2}=0,
\\[2mm]
\displaystyle 
\frac{C_{55}^{( s )} }{C_{11}^{( s )} }\nabla 
\frac{\partial W}{\partial r}+
\frac{C_{33}^{( s )} }{C_{11}^{( s )} } 
\frac{\partial ^2W}{\partial z^2}+
\frac{C_{13}^{( s )} +C_{55}^{( s )} }{C_{11}^{( s )}}\frac{\partial }{\partial z}\nabla U+
\\[2mm]
\displaystyle \hspace{3cm}
+\frac{e_{15} }{e_{33} }\nabla \frac{\partial \varphi }{\partial r}+
\frac{\partial ^2\varphi }{\partial z^2}-\Phi^{( s )}\frac{\partial ^2W}{\partial t^2}=0,
\\[2mm]
\displaystyle 
\frac{  e_{31} +e_{15} }{e_{33}}
\frac{\partial }{\partial  z}\nabla U+\frac{e_{15} }{e_{33} }\nabla 
\frac{\partial W}{\partial r}+\frac{\partial^2W}{\partial z^2}-
\\[2mm]
\displaystyle \hspace{3cm}
- \frac{C_{11}^{( 1 )} \varepsilon _{11}}{e_{33}^2}\nabla \frac{\partial \varphi }{\partial 
r}-
\frac{C_{11}^{( 1 )} \varepsilon _{33}}{e_{33}^2}\frac{\partial^2\varphi }{\partial z^2}=0,
\end{array}
\end{equation}
где 
$$\{U,W,r,z\}= \{U^\ast ,W^\ast ,r_*, z_* \}/b,\quad \varphi =\varphi^\ast e_{33}/(b C_{11}^{( 1 )} ),$$ 
$ t= t_\ast b^{-1} \sqrt{C_{11}^{(1)}/\rho^{(1)}} $ "--- время;
$U^\ast ( r_\ast ,z_\ast ,t_\ast )$, 
$W^\ast (r_\ast,z_\ast ,t_\ast )$ "--- радиальная и~аксиальная компоненты вектора перемещений; 
$\varphi^\ast \big(r_\ast ,z_\ast ,t_\ast \big)$ "--- 
потенциал электрического поля; 
$e_{mk}$, $\varepsilon _{11}$, $\varepsilon_{33}$ "--- пьезомодули и коэффициенты диэлектрической проницаемости 
электроупругого материала, $m ,k\hm =1,2, \dots, 5$; 
$\rho^{( s )}$, $C_{mk}^{( s )}$ "--- объемная плотность и модули  упругости электроупругого $(s=1)$ и упругого $(s=2)$ материалов; 
$\Phi^{( 1 )}=1$, $\Phi^{( 2)}\hm =\tfrac{C_{11}^{( 2 )} \rho ^{( 1 )}}{C_{11}^{( 1 )} 
\rho ^{( 2 )}}$, $\nabla=\frac{\partial }{\partial r}+\frac{1}{r}$;

\item механические граничные условия:
\begin{equation}
\label{shlya:eq2}
%r=0,1 \quad 
|W( 0,z,t )|<\infty, \quad |U(0,z,t)|<\infty;
\end{equation}

жесткое закрепление: 
\begin{equation}
\label{shlya:eq3}
W(1,z,t)=0,\quad U(1,z,t)=0;
\end{equation}

шарнирное закрепление: 
\begin{equation}
\label{shlya:eq4}
\begin{array}{l}
W(1,z,t)=0,
\\[2mm]
\displaystyle 
\sigma _{rr} \big|_{r=1} =\Bigl[ \frac{C_{11}^{( s )} }{C_{11}^{( 1 
)} }\frac{\partial U}{\partial r}+\frac{C_{12}^{( s )} 
}{C_{11}^{( 1 )} }\frac{U}{r}+\frac{C_{13}^{( s )} 
}{C_{11}^{( 1 )} }\frac{\partial W}{\partial z}+\frac{e_{31} 
}{e_{33} }\frac{\partial \varphi }{\partial z}\Bigr]_{r=1} =0;
\end{array}
\end{equation}

\begin{equation}
\label{shlya:eq5}
\begin{array}{l}
\displaystyle
\sigma _{zz}\big|_{z=0} =
\Bigl[ \frac{C_{13}^{( 1 )} }{C_{11}^{( 1 )} 
}\nabla U+\frac{C_{33}^{( 1 )} }{C_{11}^{( 1 )} 
}\frac{\partial W}{\partial z}+\frac{\partial \varphi }{\partial z} 
\Bigr]_{z=0}
=q,
\\[5mm]
\displaystyle
\sigma _{rz}\big|_{z=0} =
\Bigl[
\frac{C_{55}^{( 1 )} }{C_{11}^{( 1 )} 
}\Bigl( \frac{\partial W}{\partial r}+\frac{\partial U}{\partial z}\Bigr)
+\frac{e_{15} }{e_{33} }\frac{\partial \varphi }{\partial r}
\Bigr]_{z=0}
=0;
\end{array}
\end{equation}

\begin{equation}
\label{shlya:eq7}
\begin{array}{l}
\displaystyle
\sigma _{zz}\big|_{z=h} =
\Bigl[\frac{C_{13}^{( 1 )} }{C_{11}^{( 1 )} 
}\nabla U+\frac{C_{33}^{( 1 )} }{C_{11}^{( 1 )} 
}\frac{\partial W}{\partial z}+\frac{\partial \varphi }{\partial z}\Bigr]_{z=h}=0,
\\[5mm]
\displaystyle
\sigma _{rz}\big|_{z=h} =
\Bigl[
\frac{C_{55}^{( 1 )} }{C_{11}^{( 1 )} }
\Bigl(
\frac{\partial W}{\partial r}+\frac{\partial U}{\partial z}\Bigr) +\frac{e_{15} }{e_{33} }\frac{\partial \varphi }{\partial r}
\Bigr]_{z=h}
=0,
\end{array}
\end{equation}

\begin{equation}
\label{shlya:eq7}
\hspace{-3mm}
\begin{array}{c}
\displaystyle
\Bigl[\frac{C_{13}^{( 1 )} }{C_{11}^{( 1 )} }\nabla  U+\frac{C_{33}^{( 1 )} }{C_{11}^{( 1 )} 
}\frac{\partial W}{\partial z}+\frac{\partial \varphi }{\partial z}\Bigr]_{z=h_1}
=
\Bigl[\frac{C_{13}^{( 2 )} }{C_{11}^{( 1 )} }\nabla  U+\frac{C_{33}^{( 2 )} }{C_{11}^{( 1 )}  }\frac{\partial W}{\partial z}\Bigr]_{z=h_1},\\[5mm]
\displaystyle
\Bigl[ \frac{C_{55}^{( 1 )} }{C_{11}^{( 1 )} }\Bigl(\frac{\partial W}{\partial r}+\frac{\partial U}{\partial z} \Bigr)+\frac{e_{15} }{e_{33}}\frac{\partial \varphi }{\partial 
r}\Bigr]_{z=h_1 +h_2 }
=
\Bigl[
\frac{C_{55}^{( 2 )} }{C_{11}^{( 1 )} }\Bigl(\frac{\partial W}{\partial r}+\frac{\partial U}{\partial z}\Bigr)
\Bigr]_{z=h_1 +h_2 },\\[4mm]
W(z+0)=W(z-0), \quad  U(z+0)=U(z-0);
\end{array} \hspace{-3mm}
\end{equation}

\item электрические граничные условия:
\begin{equation}
\label{shlya:eq8}
\begin{array}{l}
\varphi (0,z,t)<\infty,
\\[2mm]
\displaystyle
D_{r} \big|_{r=1} =
\Bigl[
-\frac{C_{11}^{( 1 )} \varepsilon _{11} 
}{e_{33}^2 }\frac{\partial \varphi }{\partial r}+\frac{e_{15} }{e_{33} 
}\Bigl(\frac{\partial W}{\partial r}+\frac{\partial U}{\partial z}\Bigr)\Bigr]_{r=1}=0;
\end{array}
\end{equation}

\begin{itemize}

\item[а)] задача обратного пьезоэффекта 
\begin{equation}
\label{shlya:eq10}
\begin{array}{c}
\quad \varphi (r,0,t)=V( r,t )/2, \quad  \varphi (r,h,t)=- V( r,t )/2, \\[2mm]
\varphi (r,h_1,t) = \varphi (r,h_1+h_2,t) =0;
\end{array}
\end{equation}

\item[б)] задача прямого пьезоэффекта 
\begin{equation}
\hspace{-1cm}
\label{shlya:eq10}
D_z\bigl|_{z=\{0, h_1, h_1+h_2, h\} } =
\Bigl[
-\frac{C_{11}^{( 1 )} \varepsilon _{33} }{e_{33}^2 
}\frac{\partial \varphi }{\partial z}+\frac{e_{31} }{e_{33} }\nabla 
U+\frac{\partial W}{\partial z}\Bigr]_{z=\{0, h_1, h_1+h_2, h\} }=0;
\end{equation}


\end{itemize}

\item начальные условия:
\begin{equation}
\label{shlya:eq11}
\begin{aligned}
& U(r,z,0)=U_0 ( r,z), \quad \frac{\partial U}{\partial t} \Big|_{t=0} =\dot {U}_0 (r,z),\\
& W(r,z,0)=W_0 (r,z), \quad \frac{\partial W}{\partial 
t}\Big|_{t=0} =\dot {W}_0 ( r,z), \end{aligned}
\end{equation}

где $\{ U_0,\dot {U}_0 ,W_0,\dot {W}_0,h,h_1 ,h_2 \}=
\{ U_0^\ast,\dot {U}_0^\ast,W_0^\ast,\dot {W}_0^\ast 
,h^\ast ,h_1^\ast ,h_2^\ast  \}/b$; 
$\sigma _{pm}$ "--- компоненты тензора напряжений, $p,m=r,z$; 
$D_r$, $D_z$ "--- компоненты  вектора индукции электрического поля; 
$U_0^\ast$, $\dot {U}_0^\ast$, $W_0^\ast$, $\dot {W}_0^\ast$ "--- известные в~начальный момент времени перемещения и их 
скорости; $q= {q^\ast }/{C_{11}^{( 1 )}}$; $V=V^\ast {e_{33}}/(bC_{11}^{( 1 )} )$.
\end{itemize}

\smallskip

При исследовании упругой среды ($s=1)$ система \eqref{shlya:eq1} состоит только из 
уравнений движения, сформулированных относительно компонент вектора 
перемещений. 
Выражения \eqref{shlya:eq2} и первое неравенство \eqref{shlya:eq8} являются условиями 
ограниченности решения в центре пластины, равенство $D_{r} |_{r=1} =0$ \eqref{shlya:eq8} 
учитывает отсутствие электродного покрытия на цилиндрической поверхности 
пластины, а выражения \eqref{shlya:eq10} в~задаче прямого пьезоэффекта означают 
подключение электродированной поверхности пьезокерамической пластины к~измерительному прибору с~большим входным сопротивлением, что соответствует режиму <<холостого хода>> (отсутствию свободных электрических 
зарядов). 
Формулы \eqref{shlya:eq7} позволяют учесть условия совместности напряжений и~перемещений в~плоскости жесткого соединения пластин. 
Кроме этого, в задаче обратного пьезоэффекта первое равенство \eqref{shlya:eq5} приравнивается к~нулю.

\smallskip

\Section{Построение решения} Определение общего интеграла рассматриваемых 
краевых задач электроупругости выполняется путем последовательного 
применения конечных интегральных преобразований: Фурье"--~Бесселя~\cite{shlya:bib23} по 
радиальной переменной и~обобщенного преобразования~\cite{shlya:bib24} по аксиальной 
координате. 
При этом на каждом этапе предварительно выполняется процедура приведения соответствующих граничных условий к виду, позволяющему использовать алгоритм преобразования. 

Преобразование Фурье"--~Бесселя~\cite{shlya:bib23} позволяет удовлетворить только смешанные 
граничные условия. 
Для выполнения данного требования необходимо произвести следующие замены граничных условий~\eqref{shlya:eq3}, \eqref{shlya:eq4}:

\begin{itemize}
\item при жестком защемлении пластины первое равенство \eqref{shlya:eq3} заменяется условием 
наличия касательных напряжений $N_1 (z,t)$ на цилиндрической 
поверхности пластины: 
\begin{equation}
\label{shlya:eq13}
\sigma _{rz}\big|_{r=1} =
\Bigl[
\frac{C_{55}}{C_{11}}\Bigl(\frac{\partial 
W}{\partial r}+\frac{\partial U}{\partial z}\Bigr)+\frac{e_{15} }{e_{33}  }\frac{\partial \varphi }{\partial r}\Bigr] =N_1 (z,t);
\end{equation}

\item в случае шарнирного закрепления последнее соотношение \eqref{shlya:eq8} заменяется 
условием наличия потенциала электрического поля $\varphi _1 ( z,t)$ на цилиндрической поверхности при $r=1$:
\begin{equation}
\label{shlya:eq13}
\varphi ( {1,z,t} )=\varphi _1 ( {z,t} ).
\end{equation}
\end{itemize}

Кроме этого, в общем случае вводятся новые функции $u(r,z,t)$, $w(r,z,t)$, $\phi (r,z,t)$, связанные с~$U(r,z,t)$, $W(r,z,t)$, $\varphi (r,z,t)$ соотношениями
\begin{gather}
U(r,z,t)=-\frac{e_{31} }{2e_{33} }\frac{\partial \varphi _1 
(z,t)}{\partial z}r\big[ H(z-h_1)+H(z-h_1 -h_2) \big]+u(r,z,t),\\
W(r,z,t)=W_1 ( t )+A_0 C_{11} \varepsilon _{11} 
r^2N_1 (z,t)+w(r,z,t), \label{shlya:eq14}\\
\varphi (r,z,t)=A_0 e_{33} e_{15} r^2N_1 ( z,t)+\phi (r,z,t),
\quad
A_0 =\left[ 2 \left(C_{55} \varepsilon _{11} +e_{15}^2 \right) \right]^{-1}.
\end{gather}
Здесь $H({}\cdot{})$ "--- единичная функция Хэвисайда~\cite{shlya:bib25}; $W_1 (t )$, $N_1 (z,t)$, $\varphi _1 (z,t)$ "--- 
неизвестные функции, определяемые в процессе решения задачи, с помощью 
которых удается удовлетворить все условия \eqref{shlya:eq3}, \eqref{shlya:eq4}, \eqref{shlya:eq8}. 



В результате подстановки \eqref{shlya:eq14} в \eqref{shlya:eq1}--\eqref{shlya:eq13} получаем новую краевую задачу 
относительно функций $u$, $w$, $\phi$. 
При этом дифференциальные уравнения \eqref{shlya:eq1}, 
граничные условия по переменной $z$ \eqref{shlya:eq5}--\eqref{shlya:eq10} становятся неоднородными 
с~правыми частям $R_1, R_2, R_3 $, $B_1, B_2,\dots, B_{10} $, начальные условия $U_0$, $\dot {U}_0$, $W_0$, $\dot {W}_0$ следует заменить на $u_0$, $\dot {u}_0$, $w_0$, $\dot {w}_0$, а краевые условия \eqref{shlya:eq2}--\eqref{shlya:eq4} принимают следующий вид:

\begin{itemize}
\item жесткое закрепление: 
\begin{equation}
\label{shlya:eq15}
U(1,z,t)=0,\quad \cfrac{\partial W}{\partial r} \Big|_{z=1}=0, \quad \frac{\partial \phi }{\partial r} \Big|_{z=1} =0;
\end{equation}

\item шарнирное закрепление:
\begin{equation}
\label{shlya:eq16}
W(1,z,t)=0, \quad \phi (1,z,t)=0, \quad \nabla U \big|_{r=1} =0.
\end{equation} 
\end{itemize}

Последнее условие \eqref{shlya:eq16} получается при $C_{11}^{( s )} 
=C_{12}^{( s )} $. 
%Только в этом случае, в дальнейшем, дополнительные внеинтегральные члены в первом трансформированном уравнении \eqref{shlya:eq1} равны нулю.

Применим к краевой задаче относительно функций $u$, $w$, $\phi $ конечные интегральные преобразования Фурье"--~Бесселя, используя следующие трансформанты: 
\begin{gather}
\label{shlya:eq17}
u_H (j_n ,z,t)=\int _0^1 u(r,z,t)rJ_1  (j_n r )dr,\\
\bigl\{ w_H (j_n ,z,t ), \phi _H (j_n ,z,t )  \bigr\}=
\int_0^1  \bigl\{ w(r,z,t),\phi (r,z,t)  \bigr\} r J_0 (j_n r )dr
\end{gather}
и формулы обращения:
\begin{gather}
\label{shlya:eq18}
u(r,z,t)=2
\sum\limits_{n=1}^\infty 
\frac{u_H (j_n,z,t )}{S (j_n )^2} J_1 (j_n r ),\\
\bigl\{ w( r,z,t), \phi (r,z,t) \bigr\}=
2\sum\limits_{n=0}^\infty 
\frac{\bigl\{ w_H (j_n ,z,t ),\phi _H (j_n ,z,t ) \bigr\}}{S(j_n )^2} J_0 (j_n r ),
\end{gather}
где $j_n$ "--- положительные нули функций $J_1 (j_n )$;
$J_0  (j_n )$ соответственно при жестком и шарнирном закреплении 
пластины, расположенные в порядке их возрастания, $n=0,1,2,\dots, j_0 =0 $; 
$S (j_n )=J_0 (j_n )$ при жестком и $S(j_n)\hm=J_1 (j_n )$  при шарнирном закреплении.


В результате получаем новую задачу относительно трансформант Фурье"--~Бесселя 
$u_H$, $w_H$, $\phi _H$, которую решаем, используя обобщенный метод конечных 
интегральных преобразований~\cite{shlya:bib24} по аксиальной координате. 
При этом  предварительно проводится также  процедура стандартизации, связанная с~приведением неоднородных граничных условий по переменной $z$ к однородным. 
Для этой цели вводятся новые функции $U_H$, $W_H$, $\varphi _H $, связанные с~$u_H$, $w_H$, $\phi _H $ соотношениями
\begin{multline}
\label{shlya:eq19}
\bigl\{ u_H (j_n z,t ), w_H (j_n z,t ), \phi _H (j_n z,t ) \bigr\}=\\
=\bigl\{ P_1 (j_n z,t ), P_2 ( j_n z,t ), P_3 (j_n z,t ) \bigr\}+\\
+\bigl\{ U_H (j_n z,t), W_H (j_n z,t ),\varphi_H (j_n z,t ) \bigr\}.
\end{multline}

В результате подстановки \eqref{shlya:eq19} в расчетные соотношения краевой задачи 
относительно $u_H$, $w_H$, $\phi _H$ получаем новую задачу для функций $U_H$, $W_H$, $\varphi _H$. При этом граничные условия с помощью выбора функций $P_1$, $P_2$, $P_3$ становятся однородными.

Начально"=краевую задачу относительно функций $U_H$, $W_H$, $\varphi _H $ 
решаем, используя обобщенный метод конечных интегральных преобразований 
\mbox{(КИП)} \cite{shlya:bib24} с использованием неизвестной трансформанты $G(\lambda 
_{in},n,t )$ и компонент $K_1 (\lambda _{in} ,z)$, $K_2 (\lambda _{in} ,z )$, $K_3 (\lambda _{in} 
,z )$ вектор"=функции ядра преобразования ($\lambda _{in}$ "--- положительные параметры, образующие счетное множество, $i=1,2,3,\dots$)
\begin{equation}
\label{shlya:eq20}
G(\lambda _{in} ,n,t )=\int _0^h (U_H K_{1in}  +W_H K_{2in} )dz,
\end{equation}
и формул обращения
\begin{gather}
\label{shlya:eq21}
\bigl\{ U_H ,W_H ,\varphi _H \bigr\}=
\sum\limits_{i=1}^\infty 
G_{in}  \bigl\{ K_{1in} ,K_{2in} ,K_{3in} \bigr\}
\| K_{in} \|^{-2},\\
\| K_{in} \|^2=\int_0^h (K_{1in}^2 +K_{2in}^2 +K_{3in}^2 )dz.
\end{gather}

Общие решения системы дифференциальных уравнений для электроупругого $(s=1)$ и упругого $(s=2)$ слоев представлены в~работах авторов~\cite{shlya:bib18,shlya:bib26}. 

Подстановка расчетных соотношений $K_1 (\lambda_{in},z )$, $K_2 (\lambda _{in},z )$, $K_3 (\lambda _{in},z )$ в граничные условия по аксиальной координате формирует однородную систему уравнений относительно постоянных интегрирования $D_{1in}, D_{2in}, \dots,  D_{16in} $. 
Разыскивая ее нетривиальное решение, получаем трансцендентное уравнение для вычисления собственных значений $\lambda _{in} $, а также выражения для постоянных $D_{1in}, D_{2in}, \dots,  D_{16in} $. 

Окончательные выражения функций $U(r,z,t)$, $W(r,z,t)$, $\varphi (r,z,t)$ получим, применяя к трансформанте 
\eqref{shlya:eq20} последовательно формулы обращения КИП \eqref{shlya:eq21} и Фурье"--~Бесселя~\eqref{shlya:eq18}. 
В~результате с учетом \eqref{shlya:eq14}, \eqref{shlya:eq19} имеем
\begin{multline}
U(r,z,t)=-\frac{e_{31} }{2e_{33} }\frac{\partial \varphi_1 (z,t)}{\partial z}r
\bigl[H(z-h_1)+H(z-h_1 -h_2 ) \bigr]+\\
+2\sum\limits_{n=1}^\infty 
\frac{J_1 (j_n r )}{S(j_n)^2} 
\biggl(P_1 +\sum\limits_{i=1}^\infty G_{in} K_{1in} \| K_{in} \|^{-2} \biggr),~~~~~~
\end{multline}
\begin{multline}
W(r,z,t)=W_1 ( t )+A_0 C_{11} \varepsilon _{11} 
r^2N_1 (z,t)+\\
+2\sum\limits_{n=0}^\infty \frac{J_0 (j_n r )}{S(j_n )^2} 
\biggl( P_2 +\sum\limits_{i=1}^\infty G_{in} K_{2in} \| K_{in} \|^{-2} \biggr),
\label{shlya:eq22}
\end{multline}
\begin{multline}
\varphi (r,z,t)=A_0 e_{33} e_{15} r^2N_1 (z,t) +\\ 
+2\sum\limits_{n=0}^\infty 
\frac{J_0 (j_n r)}{S(j_n )^2} 
\biggl(P_3 +\sum\limits_{i=1}^\infty
G_{in} K_{3in} \| K_{in} \|^{-2}  \biggr).~~~~~~
\end{multline}

Заключительным этапом исследования является определение зависимостей $W_1 ( t)$, $N_1 (z,t )$, $\varphi _1 (z,t )$. 
При  решении задачи обратного пьезоэффекта первоначально решается задача в~случае 
действия только электрической нагрузки $V(r,t)$ (${N_1 =0}$). 

В результате определяется функция $W_1 ( t )$ из условия 
отсутствия вертикальных перемещений цилиндрической поверхности пластины при 
$z=h/2$, а также характер изменения касательных напряжений $\sigma _{rz} 
(1,z,t)$, который определяет в дальнейшем выражение для $N_1 (z,t)$. На втором этапе решения рассматривается действие функции $N_1 
(z,t)$, которая аппроксимируется параболической зависимостью 
и определяется из условия равенства нулю перемещений на цилиндрической 
поверхности пластины в случае действия $V(r,t)$ и $N_1 (z,t)$. Суммирование двух расчетов позволяет получить окончательное 
решение.

Аналогичным образом решается задача прямого пьезоэффекта. Первоначально 
рассматривается действие механической нагрузки $q(r,t)$, а на следующем 
этапе функция $\varphi _1 (r,t)$ определяется из условия равенства нулю 
радиальной компоненты вектора индукции электрического поля на цилиндрической 
поверхности при суммировании двух расчетов.

Полученные расчетные соотношения~\eqref{shlya:eq22} удовлетворяют дифференциальные 
уравнения \eqref{shlya:eq1} и краевые условия \eqref{shlya:eq2}--\eqref{shlya:eq11}, 
т.~е. представляют замкнутые решения рассматриваемых начально"=краевых задач электроупругости для 
биморфных пьезокерамических пластин.

\smallskip

\Section[n]{Заключение} В настоящей работе на основании разработанного в работах авторов \cite{shlya:bib18,shlya:bib21,shlya:bib22} алгоритма получены расчетные соотношения для решения частных задач. 
Представленная методика расчета позволяет получать замкнутые решения связанных осесимметричных нестационарных задач в пространственной постановке для круглых сплошных биморфных пластин на шарнирном и жестком закреплении по контуру.

Допущение об установившемся режиме вынужденных колебаний, используемое при 
исследовании динамических задач, справедливо только в случае, когда частоты 
внешнего гармонического воздействия существенно меньше первой частоты 
собственных колебаний.

Полученные расчетные соотношения позволяют определить 
напряженно"=деформированное состояние, спектр частот собственных колебаний и~характер распределения электрического поля в многослойных пьезокерамических  преобразователях резонансного и~нерезонансного классов в~случае действия 
нестационарной нагрузки.

%4. Численные результаты расчета позволяют проанализировать связанность полей 
%электроупругих напряжений и научно обосновать гипотезы о характере изменения 
%электрического поля при расчете составных многослойных конструкций 
%ступенчато переменной толщины с помощью прикладных теорий для тонких пластин 
%\cite{shlya:bib21,shlya:bib22}. В частности, можно отметить:
%
%\begin{itemize}
%\item для наиболее эффективного преобразования внешнего электрического 
%воздействия в механические колебания в жестко закрепленных биморфных 
%пластинах, если <<рабочей>> частотой преобразователя является первая частота 
%собственных осесимметричных колебаний отношение радиусов пьезокерамической 
%пластины и металлической подложки равно 0.7;
%
%\item кинематическая гипотеза плоских сечений выполняется при расчете биморфных 
%пластин, выполненных из материалов с разными физическими характеристиками;
%
%\item при расчете пьезокерамических тонких пластин можно использовать гипотезы о 
%линейном характере изменения потенциала и постоянном значении аксиальной 
%компоненты вектора индукции электрического поля по толщине пьезокерамической 
%пластины. При этом радиальными компонентами электрического поля можно 
%пренебречь, поскольку их численные значения существенно меньше аксиальных 
%составляющих.
%
%\end{itemize}


\AdditionalContent{Конкурирующие интересы}{Заявляем, что в отношении авторства и публикации
этой статьи конфликта интересов не имеем.} 
\AdditionalContent{Авторский вклад и ответственность}{Все  авторы  принимали  участие  в  разработке  концепции  \mbox{статьи}  и  в  написании  рукописи. Авторы несут  полную  ответственность  за  предоставление  окончательной рукописи в печать. Окончательная  версия рукописи была одобрена всеми авторами.} 
\AdditionalContent{Финансирование}{Исследование выполнялось без финансирования.}

\vspace{-3mm}


\begin{thebibliography}{99}

\Bibitem{shlya:bib1}
\paper  Bimorph and trimorph piezoelements
\yr 2011 
\inbook  Piezoceramic Sensors. Microtechnology and MEMS
\publ Springer
\publaddr Berlin, Heidelberg
\pages 179--229
\by Sharapov~V.
\crossref{10.1007/978-3-642-15311-2_6}

\Bibitem{shlya:bib2}
\paper  Devices to control and diagnose bimorph piezoelements
\yr 2013 
\inbook  Piezo-Electric Electro-Acoustic Transducers. Microtechnology and MEMS
\publ Springer
\publaddr Cham
\pages 191--212
\by Sharapov~V., Sotula~Z., Kunickaya~L.
\crossref{10.1007/978-3-319-01198-1_11}


\Bibitem{shlya:bib3}
\paper  Miniature piezo composite bimorph actuator for elevated temperature operation
\yr 2007 
\bookvol 10
\volinfo Mechanics of Solids and Structures, Parts A and B
\by Seeley~C.~E., Delgado~E., Kunzman~J., Bellamy~D.
\crossref{10.1115/imece2007-44088}
\papernumber IMECE2007-44088
\pages 405--415
\inbook International Mechanical Engineering Congress and Exposition
\procinfo Seattle, Washington, USA, November 11–15, 2007



\RBibitem{shlya:bib4}
\by Гринченко~В.~Т., Улитко~А.~Ф., Шульга~Н.~А.
\book Механика связанных полей в элементах конструкций
\publaddr Киев
\publ Наук. думка
\yr 1989
\totalpages 279

\Bibitem{shlya:bib5}
\by Smits~J.~G., Dalke~S.~I., Cooney~T.~K.
\jour Sensors and Actuators A: Physical
\vol 28
\issue 1
\yr 1991
\pages 41--61
\paper The constituent equations of piezoelectric bimorphs
\crossref{10.1016/0924-4247(91)80007-C}

\RBibitem{shlya:bib6}
\by Ватульян~А.~О., Рынкова~А.~А.
\paper Об одной модели изгибных колебаний пьезоэлектрических биморфов с разрезными электродами и ее приложениях
\jour Изв. РАН. МТТ
\yr 2007
\issue 4
\pages 114--122

\Bibitem{shlya:bib7}
\by Tsaplev~V., Konovalov~R., Abbakumov~K.
\paper Disk bimorph-type piezoelectric  energy harvester
\jour  Journal of Power and Energy Engineering
\yr 2015
\vol 3
\issue 4
\pages 63--68
\crossref{10.4236/jpee.2015.34010}

\RBibitem{shlya:bib8}
\by Цаплев~В.~М., Коновалов~Р.~С.
\paper Частотные зависимости констант высших порядков пьезоэлектрической керамики
\jour Дефектоскопия
\yr  2017
\issue 7
\pages 14–22
%https://link.springer.com/article/10.1134%2FS1061830917070087
   

\RBibitem{shlya:bib9}
\by Янчевский~И.~В.
\paper Минимизация прогибов электроупругой биморфной пластины при импульсном нагружении
\jour Проблемы вычислительной механики и прочности конструкций
\yr 2011
\issue 16
\publaddr 303--313
%https://pommk.dp.ua/index.php/journal/article/view/263/306

\RBibitem{shlya:bib10}
\by Шляхин~Д.~А.
\paper Вынужденные осесимметричные колебания пьезокерамической тонкой биморфной пластины
\jour Изв. РАН. МТТ
\yr 2013
\issue 2
\pages 77--85 

\Bibitem{shlya:bib11}
\by Rashidifar~M.~A., Rashidifar~A.~A.
\paper Vibrations analysis of circular plate with piezoelectric actuator using thin plate theory and Bessel function
\jour American Journal of Engineering, Technology and Society
\yr 2015
\issue 6
\vol 2
\pages 140--156



\Bibitem{shlya:bib12}
\by Jam~J.~E., Khosravi~M., Namdara~N.
\paper  An exact solution of mechanical buckling for functionally graded material bimorph circular 
plates
\jour Metall. Mater. Eng.
\yr 2013
\vol 19
\issue 1
\pages 45--63

\Bibitem{shlya:bib14}
\by Jandaghiana~A.~A., Jafarib~A.~A., Rahmania~O.
\paper Vibrational response of functionally graded circular plate integrated with piezoelectric layers: An 
exact solution
\jour Engineering Solid Mechanics
\yr 2014
\issue 2
\vol 2 
\pages 120--130
\crossref{10.5267/j.esm.2014.1.004}

\Bibitem{shlya:bib15}
\by  Jurenas~V., Bansevicius~R., Navickaite~S.
\paper Piezoelectric bimorphs for laser shutter systems: optimization of dynamic characteristics
\jour Mechanika
\yr 2010 
\issue 5(85)
\pages 44--47 


\RBibitem{shlya:bib17}
\by Сеницкий~Ю.~Э., Шляхин~Д.~А.
\paper Нестационарная осесимметричная задача  электроупругости для толстой круглой анизотропной пьезокерамической пластины 
\jour Изв. РАН. МТТ
\yr 1999
\issue 1
\pages 78--87

\RBibitem{shlya:bib18}
\by Шляхин~Д.~А.
\paper Вынужденные осесимметричные колебания толстой круглой жестко 
закрепленной пьезокерамической пластины
\jour Изв. РАН. МТТ
\yr 2014
\issue 4
\pages 90--100

\RBibitem{shlya:bib19}
\by Шульга~Н.~А., Болкисев~А.~М.
\book Колебания пьезоэлектрических тел
\publaddr Киев
\publ Наук. думка
\yr 1990
\totalpages 228

\Bibitem{shlya:bib21}
\by Shlyakhin~D.~A., Kazakova~O.~V.
\paper A dynamic axially symmetric goal and its  extended solution for a fixed rigid circular multi-layer plate
\jour Procedia Engineering
\yr 2016 
\vol 153
\pages 662--666
\crossref{10.1016/j.proeng.2016.08.219}

\RBibitem{shlya:bib22}
\by Шляхин~Д.~А.
\paper Динамическая осесимметричная задача прямого пьезоэффекта для  круглой биморфной пластины
\jour Вестник Пермского национального исследовательского политехнического университета. Механика
\issue 1
\yr 2017
\pages 164--180
\crossref{10.15593/perm.mech/2017.1.10}

\Bibitem{shlya:bib23}
\book Fourier Transforms
\by Sneddon~I.~N. 
\zmath{0038.26801}
\publaddr New York
\publ McGraw-Hill Book Company, Inc.
\yr 1950

\RBibitem{shlya:bib24}
\by Сеницкий~Ю.~Э.
\book Исследование упругого деформирования элементов конструкций при динамических воздействиях методом конечных интегральных преобразований
\publaddr Саратов
\publ Сарат. ун-т
\yr 1985
\totalpages 174

\RBibitem{shlya:bib25}
\by Владимиров~В.~С.
\book Обобщенные функции в математической физике
\publaddr М.
\publ Наука
\yr 1978
\totalpages 318


\RBibitem{shlya:bib26}
\by Шляхин~Д.~А.
\paper Вынужденные осесимметричные изгибные колебания толстой круглой жестко закрепленной пластины
\jour Вестн. СамГУ. Естественнонаучн. сер.
\yr 2011
\issue 8(89)
\pages 142--152
\mathnet{http://mi.mathnet.ru/vsgu133}


\end{thebibliography}


\newpage


\makeentitle

\vspace{-5mm}

\AdditionalContent{Competing interests}{We declare that we have no conflicts of interests with the authorship and publication of this article.} 
\AdditionalContent{Authors' contributions and responsibilities}{Each author has participated in the article concept development and in the manuscript writing. The authors are absolutely responsible for submitting the final manuscript in print. Each author has approved the final version of manuscript.} 

\AdditionalContent{Funding}{The research has not had any funding.}



\begin{thebibliography}{99}

\Bibitem{shlya:bib1en}
\paper  Bimorph and trimorph piezoelements
\yr 2011 
\inbook  Piezoceramic Sensors. Microtechnology and MEMS
\publ Springer
\publaddr Berlin, Heidelberg
\pages 179--229
\by Sharapov~V.
\crossref{10.1007/978-3-642-15311-2_6}

\Bibitem{shlya:bib2en}
\paper  Devices to control and diagnose bimorph piezoelements
\yr 2013 
\inbook  Piezo-Electric Electro-Acoustic Transducers. Microtechnology and MEMS
\publ Springer
\publaddr Cham
\pages 191--212
\by Sharapov~V., Sotula~Z., Kunickaya~L.
\crossref{10.1007/978-3-319-01198-1_11}


\Bibitem{shlya:bib3en}
\paper  Miniature piezo composite bimorph actuator for elevated temperature operation
\yr 2007 
\bookvol 10
\volinfo Mechanics of Solids and Structures, Parts A and B
\by Seeley~C.~E., Delgado~E., Kunzman~J., Bellamy~D.
\crossref{10.1115/imece2007-44088}
\papernumber IMECE2007-44088
\pages 405--415
\inbook International Mechanical Engineering Congress and Exposition
\procinfo Seattle, Washington, USA, November 11–15, 2007


\Bibitem{shlya:bib4en}
\lang In Russian
\by Grinchenko~V.~T., Ulitko~A.~F., Shul'ga~N.~A.
\book Mekhanika sviazannykh polei v elementakh konstruktsii \rm [Mechanics of coupled fields in structural elements]
\publaddr Kiev
\publ Nauk. dumka
\yr 1989
\totalpages 279

\Bibitem{shlya:bib5:en}
\by Smits~J.~G., Dalke~S.~I., Cooney~T.~K.
\jour Sensors and Actuators A: Physical
\vol 28
\issue 1
\yr 1991
\pages 41--61
\paper The constituent equations of piezoelectric bimorphs
\crossref{10.1016/0924-4247(91)80007-C}

\Bibitem{shlya:bib6:en}
\jour Mech. Solids 
\yr 2007
\vol 42
\crossref{10.3103/S0025654407040127}
\issue 4
\pages 595–602
\paper A model of bending vibrations of piezoelectric bimorphs with split electrodes and its applications
\by Vatul'yan~A.~O., Rynkova~A.~A



\Bibitem{shlya:bib7en}
\by Tsaplev~V., Konovalov~R., Abbakumov~K.
\paper Disk bimorph-type piezoelectric  energy harvester
\jour  Journal of Power and Energy Engineering
\yr 2015
\vol 3
\issue 4
\pages 63--68
\crossref{10.4236/jpee.2015.34010}

\Bibitem{shlya:bib8:en}
\by Tsaplev~V.~M., Konovalov~R.~S.
\jour  Russ. J. Nondestruct. Test 
\yr 2017
\vol 53
\crossref{10.1134/S1061830917070087}
\issue 7
\pages 485–492 
\paper Frequency dependences of higher-order constants of piezoelectric ceramics

   

\Bibitem{shlya:bib9en}
\lang In Russian
\by Yanchevskiy~I.~V.
\paper Minimizing deflections of round electroelastic bimorph plate under impulsive loading
\jour  Problems of computational mechanics and strength of structures 
\yr 2011
\issue 16
\publaddr 303--313
%https://pommk.dp.ua/index.php/journal/article/view/263/306

\Bibitem{shlya:bib10:en}
\by Shlyakhin~D.~A.
\jour  Mech. Solids 
\yr 2013
\vol 48
\crossref{10.3103/S002565441302009X}
\issue 2
\pages 178–185 
\paper Forced nonstationary axisymmetric vibrations of a piezoceramic thin bimorph plate

\Bibitem{shlya:bib11en}
\by Rashidifar~M.~A., Rashidifar~A.~A.
\paper Vibrations analysis of circular plate with piezoelectric actuator using thin plate theory and Bessel function
\jour American Journal of Engineering, Technology and Society
\yr 2015
\issue 6
\vol 2
\pages 140--156



\Bibitem{shlya:bib12en}
\by Jam~J.~E., Khosravi~M., Namdara~N.
\paper  An exact solution of mechanical buckling for functionally graded material bimorph circular 
plates
\jour Metall. Mater. Eng.
\yr 2013
\vol 19
\issue 1
\pages 45--63

\Bibitem{shlya:bib14en}
\by Jandaghiana~A.~A., Jafarib~A.~A., Rahmania~O.
\paper Vibrational response of functionally graded circular plate integrated with piezoelectric layers: An 
exact solution
\jour Engineering Solid Mechanics
\yr 2014
\issue 2
\vol 2 
\pages 120--130
\crossref{10.5267/j.esm.2014.1.004}

\Bibitem{shlya:bib15en}
\by  Jurenas~V., Bansevicius~R., Navickaite~S.
\paper Piezoelectric bimorphs for laser shutter systems: optimization of dynamic characteristics
\jour Mechanika
\yr 2010 
\issue 5(85)
\pages 44--47 


\Bibitem{shlya:bib17:en}
\by Senitskii~Yu.~E., Shlyakhin~D.~A.
\paper The nonstationary axisymmetric problem of electroelasticity for a thick circular anisotropic piezoceramic plate
\jour Mech. Solids 
\vol 34
\issue 1
\pages 66--74 
\yr 1999

\Bibitem{shlya:bib18:en}
\by Shlyakhin~D.~A.
\jour  Mech. Solids 
\yr 2014
\vol 49
\crossref{10.3103/S0025654414040086}
\issue 4
\pages 435–444 
\paper Forced axisymmetric vibrations of a thick circular rigidly fixed piezoceramic plate

\Bibitem{shlya:bib19:en}
\lang In Russian
\by Shul'ga~N.~A., Bolkisev~A.~M.
\book Kolebaniia p'ezoelektricheskikh tel \rm [Fluctuations of Piezoelectric Bodies]
\publaddr Kiev
\publ Nauk. dumka
\yr 1990
\totalpages 228


\Bibitem{shlya:bib21en}
\by Shlyakhin~D.~A., Kazakova~O.~V.
\paper A dynamic axially symmetric goal and its  extended solution for a fixed rigid circular multi-layer plate
\jour Procedia Engineering
\yr 2016 
\vol 153
\pages 662--666
\crossref{10.1016/j.proeng.2016.08.219}

\Bibitem{shlya:bib22en}
\lang In Russian
\by Shlyakhin~D.~A.
\paper Dynamic axisymmetric problem direct piezoeffect for round bimorph plate
\jour PNRPU Mechanics Bulletin
\issue 1
\yr 2017
\pages 164--180
\crossref{10.15593/perm.mech/2017.1.10}

\Bibitem{shlya:bib23en}
\book Fourier Transforms
\by Sneddon~I.~N. 
\zmath{0038.26801}
\publaddr New York
\publ McGraw-Hill Book Company, Inc.
\yr 1950

\Bibitem{shlya:bib24en}
\lang In Russian
\by Senitskii~Yu.~E.
\book Issledovanie uprugogo deformirovaniia elementov konstruktsii pri dinamicheskikh vozdeistviiakh metodom konechnykh integral'nykh preobrazovanii \rm [Investigation of the Elastic Deformation of Construction Elements under Dynamic Actions by the Method of Finite Integral Transforms]
\publaddr Saratov
\publ Saratov Univ.
\yr 1985
\totalpages 174



\Bibitem{shlya:bib25en}
\lang In Russian
\by Vladimirov~V.~S.
\book Obobshchennye funktsii v matematicheskoi fizike \rm [Generalized Functions in Mathematical Physics]
\publaddr Moscow
\publ Nauka
\yr 1978
\totalpages 318



\Bibitem{shlya:bib26en}
\lang In Russian
\by Shlyakhin~D.~A.
\paper The compelled axisymmetric bending fluctuations of thick round rigid plate
\jour Vestnik Samarskogo Gosudarstvennogo Universiteta. Estestvenno-Nauchnaya Seriya
\yr 2011
\issue 8(89)
\pages 142--152


\end{thebibliography}

\renewcommand{\baselinestretch}{0.9}

%\end{document}

